% ================================================================
% SECTION 15: TIME INTEGRATION
% ================================================================
\section{Time Integration}
\label{sec:time_integration}

\begin{center}
\code{bionetflux.time\_integration.newton\_solver} \quad(445 lines) \\
\code{bionetflux.time\_integration.time\_stepper} \quad(462 lines)
\end{center}

Time integration uses an implicit Euler scheme solved at each step by a
Newton method.  The module provides two levels of abstraction: a low-level
\code{NewtonSolver} and a higher-level \code{TimeStepper} /
\code{AdaptiveTimeStepper}.

% ==================================================================
\subsection{NewtonSolver}

\begin{lstlisting}[language=Python, caption={NewtonSolver constructor}]
class NewtonSolver:
    def __init__(self, tolerance: float = 1e-10,
                 max_iterations: int = 20,
                 verbose: bool = False):
\end{lstlisting}

\subsubsection{NewtonResult dataclass}

\begin{lstlisting}[language=Python, caption={NewtonResult fields}]
@dataclass
class NewtonResult:
    converged: bool
    iterations: int
    final_solution: np.ndarray
    final_residual_norm: float
    residual_history: List[float]
    step_norms: List[float]
    jacobian_condition: Optional[float] = None
    computation_time: float = 0.0
\end{lstlisting}

\subsubsection{Solver Variants}

Three methods share the same interface:

\begin{lstlisting}[language=Python, caption={Common signature}]
def solve(self, initial_guess, global_assembler, forcing_terms,
          static_condensations, current_time, ...) -> NewtonResult:
\end{lstlisting}

\begin{longtable}{p{7.5cm}p{7cm}}
\toprule
\textbf{Method} & \textbf{Strategy} \\
\midrule
\code{solve} & Standard Newton: $J\,\delta x = -R$; update $x \mathrel{+}= \delta x$. Uses \code{np.linalg.solve}. \\
\code{solve\_with\_line\_search(alpha\_init=1.0, alpha\_min=1e-4)} & Backtracking line search: halves $\alpha$ until $\|R(x + \alpha\,\delta x)\| < \|R(x)\|$. \\
\code{solve\_with\_damping(damping\_factor=0.8)} & Fixed damping: $x \mathrel{+}= \omega\,\delta x$ with constant $\omega \in (0,1]$. \\
\bottomrule
\end{longtable}

All variants monitor the Jacobian condition number and record the full residual history.

% ==================================================================
\subsection{TimeStepResult dataclass}

Defined inline in \code{time\_stepper.py}:

\begin{lstlisting}[language=Python, caption={TimeStepResult fields}]
@dataclass
class TimeStepResult:
    converged: bool
    iterations: int
    final_residual_norm: float
    updated_solution: np.ndarray
    updated_bulk_data: List          # List[BulkData]
    computation_time: float
    # Optional:
    residual_history: Optional[List[float]] = None
    jacobian_condition: Optional[float] = None
    newton_step_norms: Optional[List[float]] = None
    flux_data: Optional[List] = None
\end{lstlisting}

\code{summary()} returns a multi-line human-readable report.

% ==================================================================
\subsection{TimeStepper}

\begin{lstlisting}[language=Python, caption={TimeStepper constructor}]
class TimeStepper:
    def __init__(self, setup: SolverSetup,
                 newton_solver: Optional[NewtonSolver] = None,
                 verbose: bool = True):
\end{lstlisting}

\noindent Caches \code{setup.bulk\_data\_manager}, \code{setup.global\_assembler},
\code{setup.static\_condensations}, \code{setup.problems}, and the spatial
discretization list.

\subsubsection{initialize\_solution}

\begin{lstlisting}[language=Python, caption={Initial condition setup}]
def initialize_solution(self) -> Tuple[np.ndarray, List[BulkData]]:
    """Evaluate ICs, assemble global vector, initialize bulk data."""
\end{lstlisting}

\subsubsection{advance\_time\_step}

Single-step implicit Euler:

\begin{lstlisting}[language=Python, caption={One time step}]
def advance_time_step(self, current_solution, current_bulk_data,
                      current_time, dt) -> TimeStepResult:
\end{lstlisting}

\noindent Internally:
\begin{enumerate}
  \item \code{bulk\_manager.compute\_source\_terms(time=t+dt)}.
  \item For each domain \code{static\_condensation.assemble\_forcing\_term(previous\_bulk\_solution,\\ external\_force)}.
  \item \code{newton\_solver.solve(\ldots)} 
  \item If converged: \code{global\_assembler.bulk\_by\_static\_condensation} to recover bulk data.
\end{enumerate}

\subsubsection{advance\_multiple\_steps}

\begin{lstlisting}[language=Python, caption={Multi-step loop}]
def advance_multiple_steps(self, initial_solution, initial_bulk_data,
                           start_time, dt, n_steps,
                           stop_on_failure=True) -> List[TimeStepResult]:
\end{lstlisting}

Iterates \code{advance\_time\_step} and accumulates \code{TimeStepResult} objects.

% ==================================================================
\subsection{AdaptiveTimeStepper}

\begin{lstlisting}[language=Python, caption={AdaptiveTimeStepper constructor}]
class AdaptiveTimeStepper(TimeStepper):
    def __init__(self, setup, newton_solver=None, verbose=True,
                 dt_min=None, dt_max=None, safety_factor=0.8):
\end{lstlisting}

Defaults: \code{dt\_min = dt\_config * 1e-4}, \code{dt\_max = dt\_config} (from the TOML configuration).

\subsubsection{advance\_time\_step\_adaptive}

\begin{lstlisting}[language=Python, caption={Adaptive step}]
def advance_time_step_adaptive(self, current_solution, current_bulk_data,
                               current_time, dt_suggested
                               ) -> Tuple[TimeStepResult, float]:
    """Returns (result, dt_next)."""
\end{lstlisting}

\noindent Adaptation rules:
\begin{itemize}
  \item \textbf{Converged in $\le 8$ iterations}: increase $\Delta t \times 1.2$ (capped at \code{dt\_max}).
  \item \textbf{Converged in 9--14 iterations}: keep $\Delta t$.
  \item \textbf{Converged in $> 14$ iterations}: decrease $\Delta t \times 0.8$.
  \item \textbf{Failed}: halve $\Delta t$ and retry (up to 5 retries).
\end{itemize}

Before each attempt the method updates \code{GlobalDiscretization.dt} and
calls \code{sc.update\_dt()} on every static condensation to rebuild the
$\Delta t$-dependent element matrices.
